\subsection{可对角化自同态}

定义5. --- 设 E 是交换域 K 上的有限维向量空间，并设 F 是一个由 E 的自同态组成的集合。称 F 关于 E 的一个基 \((e_{i})\) 是对角的，如果每个 \(u \in F\) 关于 \((e_{i})\) 的矩阵都是对角的。称集合 F 是可对角化的，如果存在 E 的一个基，使 F 关于该基是对角的。

这个定义特别适用于 \(\mathfrak{F}\) 只含一个元素 \(u\) 的情形；这时我们称 \(u\) 是对角的（可对角化的）。还要注意，\(\mathfrak{F}\) 关于基 \((e_i)\) 是对角的，当且仅当诸 \((e_i)\) 是 \(\mathfrak{F}\) 的所有元素公共的特征向量；由此可得，\(\mathfrak{F}\) 可对角化当且仅当 \(E\) 由 \(\mathfrak{F}\) 的所有元素公共的特征向量生成。

设 A 是 \(\operatorname{End}_{\mathbf{K}}(\mathbf{E})\) 的一个包含 \(Id_{E}\) 的子代数。则 A 可对角化当且仅当它同构于某个代数 \(K^{r}\)（换言之，在 V，第 28 页，定义 1 的意义下是可对角化的）；事实上，若 A 同构于 \(K^{r}\)，则由 V，第 29 页，命题 1，A 是可对角化的；反之，若 A 可对角化，则 A 同构于对角矩阵代数的一个子代数，而对角矩阵代数作为代数同构于 \(K^{n}\)，\(n = \dim(\mathbf{E})\)，因此 A 同构于某个代数 \(K^{r}\)（V，第 30 页，命题 3）。

命题12. --- 设 E 是交换域 K 上的有限维向量空间，且 u 是 E 的一个自同态。则以下条件等价：

\begin{enumerate}
\def\labelenumi{(\roman{enumi})}
\item
  \(u\) 是可对角化的。
\item
  E 是 \(u\) 的诸特征子空间的直和。
\item
  \(u\) 的最小多项式的所有根都位于 \(K\)，且这些根都是单根。
\end{enumerate}

此外，若这些条件满足，则 E 中在 u 下封闭的每个子空间都是其与 u 的特征子空间的交的直和。

(i) 与 (ii) 的等价性由前述注记及 VII，第 31 页，命题 2 得出。设 u 可对角化，并令 \((\alpha_{i})\) 为其特征值族，\((\mathbf{V}_{i})\) 为对应的特征子空间族；由于 u 在 \(V_{i}\) 上的限制是由 \(\alpha_{i}\) 定义的位似变换，它零化多项式 \(X - \alpha_{i}\)；由此可知 u 零化多项式 \(\prod_{i}(X - \alpha_{i})\)，后者因此是 u 的最小多项式的倍式，从而与之重合，这就证明了 (iii)。反之，若 (iii) 成立，则存在 E 的一组基，使 u 的矩阵 U 是若尔当矩阵 \(U_{m,\alpha}\) 的一个对角分块（VII，第 35 页，命题 8）；于是由命题 9，整数 m 都等于 1，因此 U 是对角的。最后，末一个断言由如下事实得出：若 u 可对角化，则其诸特征子空间是 \(E_{u}\) 的准素分量，并由 VII，第 9 页，推论 1 得出。

推论。--- 若 u 的特征多项式的所有根都在 K 中，且它们都是单根，则 u 是可对角化的。

事实上，最小多项式整除特征多项式。

命题13. --- 设 E 是交换域 K 上的有限维向量空间，设 F 是一个由 E 的自同态组成的集合，并设 A 为 \(\mathrm{End}_{\mathbf{K}}(\mathbf{E})\) 中由 F 和 \(\mathrm{Id}_{\mathbf{E}}\) 生成的子代数。则以下条件等价：

\begin{enumerate}
\def\labelenumi{(\roman{enumi})}
\item
  \(\mathfrak{F}\) 是可对角化的。
\item
  K-代数 \(\mathbf{A}\) 是可对角化的。
\item
  \(\mathfrak{F}\) 的元素都是可对角化的，并且两两交换。
\end{enumerate}

如果 \((e_{i})\) 是 E 的一组基，F 关于它是对角的，则 A 包含在关于这组基为对角的自同态代数中，因而也是可对角化的；若 A 可对角化，则同样的论证表明 F 可对角化。这证明了 (i) 与 (ii) 的等价。由于任意两个对角矩阵可交换，我们有 \((\mathrm{i}) \Rightarrow (\mathrm{iii})\)，剩下要证其逆。设 F 的元素都可对角化且彼此交换。我们将利用以下引理：

引理3. --- 设 g 和 h 是向量空间 E 的两个可交换的自同态。则 g 的每个特征子空间在 h 下封闭。

事实上，如果 \(\mathbf{W}_{\lambda}\) 是 \(g\) 的对应于特征值 \(\lambda\) 的特征子空间，则对所有 \(x \in \mathbf{W}_{\lambda}\)，我们有

\[
g h. x = h g. x = h. \lambda x = \lambda h. x,
\]

这表明 \(h \cdot x \in \mathbf{W}_{\lambda}\)。

现在我们回到命题 13 的证明。在把 E 分解为非零子空间（每个都在 \(\mathfrak{F}\) 的所有元素下封闭）的直和的所有分解中，选取分量数最多的一个（E 的维数是这个数的上界），记为 \(E = \sum_{i\in I}E_i\)。设 \(u\in \mathfrak{F}\)，并令 \(E = \sum_{\alpha}V_{\alpha}\) 为

把 E 分解为 u 的特征子空间的直和的分解。由引理 3，每个 \(V_{\alpha}\) 在 F 下封闭，因而每个 \(V_{\alpha} \cap E_{i}\) 也在 F 下封闭；由命题 12，每个 \(E_{i}\) 是诸 \(V_{\alpha} \cap E_{i}\) 的直和。于是 \(E_{i}\) 的上述选取迫使每个 \(E_{i}\) 包含在某个 \(V_{\alpha}\) 之中；从而 u 在每个 \(E_{i}\) 上的限制是一个位似变换。由于这对 F 的所有元素都成立，由此可知 F 是可对角化的。

推论。--- E 的两个可交换的可对角化自同态的和与复合都是可对角化的。

